%*************************************
% ภาคผนวก ซ.
%*************************************
\vspace{-0.5in}

\appendixChapterTitle{ค}{หน้าที่ของแต่ละรีโพซิทอรีและแนวทางการตรวจสอบปัญหาเบื้องต้น}
\section{หน้าที่ของแต่ละรีโพซิทอรีในระบบ}
\hspace*{1.3em}
การแยกระบบออกเป็นหลายรีโพซิทอรีช่วยลด coupling และทำให้ทีมพัฒนาสามารถดูแลแต่ละส่วนของแพลตฟอร์มได้อย่างเป็นอิสระมากขึ้น อย่างไรก็ตาม ผู้พัฒนาต้องเข้าใจความรับผิดชอบของแต่ละรีโพซิทอรีอย่างชัดเจนเพื่อให้ตรวจสอบปัญหาได้ถูกจุด

\begin{mycustomenum2}
    \item \textbf{Midori} รับผิดชอบส่วนติดต่อผู้ใช้ การเรียก API ของระบบ และประสบการณ์ผู้ใช้บนเว็บ
    \item \textbf{Momoi} รับผิดชอบตรรกะทางธุรกิจ การเชื่อมต่อฐานข้อมูล การเปิดเผย OpenAPI การยืนยันตัวตนผ่าน `/api/auth` และการส่งงานเข้าสู่คิว
    \item \textbf{Yuzu} รับผิดชอบการประมวลผลงานเบื้องหลัง เช่น provision, deprovision และการเปลี่ยนสถานะเครื่องเสมือน
    \item \textbf{Satsuki} รับผิดชอบการสร้างไฟล์ reverse proxy configuration และ `authorized\_keys` จากข้อมูลในฐานข้อมูล โดยติดตั้งบนเครื่อง Nginx ภายนอกที่เป็น reverse proxy หน้าด่านของระบบ
    \item \textbf{Yuuka} รับผิดชอบ manifest, Docker Compose, ตัวอย่าง environment และการตั้งค่า observability ของระบบ
\end{mycustomenum2}

\section{ส่วนประกอบสนับสนุนที่เกี่ยวข้องกับหลายรีโพซิทอรี}
\hspace*{1.3em}
แม้ว่าบริการหลักเชิงธุรกิจจะมีเพียงสามตัว แต่การทำงานจริงยังขึ้นกับบริการสนับสนุนหลายตัวที่เชื่อมโยงกันระหว่างรีโพซิทอรี เช่น Satsuki สำหรับ reverse proxy/SSH automation บนเครื่อง Nginx ภายนอก, PostgreSQL สำหรับข้อมูลเชิงธุรกรรม Redis และ BullMQ สำหรับคิวงาน RustFS สำหรับพื้นที่จัดเก็บไฟล์ และ Jaeger, Prometheus, Grafana, Loki สำหรับการติดตามสถานะของระบบ

\section{แนวทางการตรวจสอบปัญหาเบื้องต้น}
\hspace*{1.3em}
เมื่อเกิดปัญหาระหว่างการพัฒนา ควรเริ่มจากการแยกก่อนว่าปัญหาเกิดในรีโพซิทอรีใด หรือเกิดจากการเชื่อมต่อข้ามรีโพซิทอรี ดังแนวทางต่อไปนี้

\begin{mycustomenum2}
    \item \textbf{เปิดหน้าเว็บไม่ได้หรือข้อมูลไม่แสดง} ให้ตรวจสอบ Midori ก่อนว่าเริ่มต้นได้สมบูรณ์ และค่า endpoint ของ Momoi ถูกต้อง
    \item \textbf{เข้าสู่ระบบไม่ได้หรือ session ผิดปกติ} ให้ตรวจสอบ Momoi, การตั้งค่า Google OAuth, cookie domain/path และการเชื่อมต่อฐานข้อมูลของข้อมูลผู้ใช้
    \item \textbf{API ตอบกลับผิดพลาด} ให้ตรวจสอบ log ของ Momoi รวมถึงสถานะของ PostgreSQL, Redis และค่าที่อยู่ในไฟล์ environment
    \item \textbf{งาน provision หรือเปลี่ยนสถานะ VM ไม่ทำงาน} ให้ตรวจสอบ Yuzu, Redis queue และการเชื่อมต่อไปยัง Proxmox VE
    \item \textbf{reverse proxy หรือการเข้าใช้งานผ่าน SSH มีปัญหา} ให้ตรวจสอบ Satsuki, trigger/notification จาก PostgreSQL, ไฟล์ configuration ที่ถูกสร้างบนเครื่อง Nginx ภายนอก และสถานะการ reload ของบริการ proxy
    \item \textbf{อัปโหลดไฟล์หรือเรียกไฟล์ไม่ได้} ให้ตรวจสอบการตั้งค่า RustFS/S3-compatible storage ทั้งฝั่ง Momoi และบริการประกอบใน Yuuka
    \item \textbf{ข้อมูล observability ไม่ขึ้น dashboard} ให้ตรวจสอบการ export metrics และ traces ของแต่ละบริการ รวมถึงสถานะของ Prometheus, Loki และ Jaeger
\end{mycustomenum2}

\section{ตัวอย่างรายการตรวจสอบอย่างเป็นลำดับ}
\hspace*{1.3em}
ในทางปฏิบัติ การไล่ตรวจสอบจากชั้นนอกเข้าสู่ชั้นในจะช่วยลดเวลาในการหาสาเหตุของปัญหาได้ โดยอาจเริ่มจาก browser, frontend, API, queue worker, database และ infrastructure ตามลำดับ ดังนี้

\begin{mycustomenum2}
    \item ตรวจสอบว่า Midori ทำงานที่ localhost:3000
    \item ตรวจสอบว่า endpoint ของ Momoi ตรงกับค่าที่ตั้งในไฟล์ environment ของ Midori
    \item ตรวจสอบ docker compose services ในรีโพซิทอรี Yuuka
    \item ตรวจสอบ Redis และ PostgreSQL
    \item ตรวจสอบ log ของ Yuzu และสถานะ queue
    \item ตรวจสอบ log ของ Satsuki และผลลัพธ์ของไฟล์ reverse proxy/authorized\_keys บนเครื่อง Nginx ภายนอกหรือ bastion ปลายทาง
    \item ตรวจสอบ Jaeger, Prometheus และ Grafana
\end{mycustomenum2}

\section{ข้อควรระวังเรื่องความสอดคล้องระหว่างรีโพซิทอรี}
\hspace*{1.3em}
เนื่องจากแต่ละรีโพซิทอรีพัฒนาแยกจากกัน จึงควรระมัดระวังความสอดคล้องของ schema ฐานข้อมูล, ตัวแปร environment, เส้นทาง API, queue name และข้อมูลอ้างอิงที่ใช้ร่วมกัน เพราะหากค่าหนึ่งค่าผิดไปเพียงจุดเดียว อาจทำให้เกิดปัญหาต่อเนื่องข้ามหลายบริการได้

\section{เอกสารอ้างอิงที่ควรใช้ประกอบ}
\hspace*{1.3em}
ในการทำงานจริง ผู้พัฒนาควรอ่าน README และเอกสารในแต่ละรีโพซิทอรีร่วมกัน โดยเฉพาะเอกสารใน Yuuka ที่สรุปตัวอย่าง environment และการติดตั้งเครื่องมือประกอบ ส่วน Midori, Momoi, Yuzu และ Satsuki ควรใช้ร่วมกับเอกสาร API, queue, authentication flow, reverse proxy และ observability เพื่อให้เห็นภาพรวมของระบบครบถ้วน

\begin{mycustomenum2}
    \item \textbf{Port ถูกใช้งานแล้ว}
    \begin{mycustomenum2}
        \item ตรวจสอบ Port ที่ใช้งานด้วยคำสั่ง: lsof -i :3000 (สำหรับ macOS/Linux)
        \item หยุดกระบวนการที่ใช้ Port หรือเปลี่ยน Port ในไฟล์ .env
    \end{mycustomenum2}

    \item \textbf{Docker Services ไม่สามารถเริ่มได้}
    \begin{mycustomenum2}
        \item ตรวจสอบสถานะด้วยคำสั่ง: docker-compose ps
        \item ดู Log ด้วยคำสั่ง: docker-compose logs [service\_name]
        \item รีสตาร์ท Services ด้วยคำสั่ง: docker-compose restart
    \end{mycustomenum2}

    \item \textbf{Database Connection Error}
    \begin{mycustomenum2}
        \item ตรวจสอบว่า PostgreSQL Container ทำงานอยู่
        \item ตรวจสอบการตั้งค่าใน Environment Variables
        \item รันคำสั่ง Database Migration (หากจำเป็น)
    \end{mycustomenum2}

    \item \textbf{Frontend ไม่สามารถเชื่อมต่อ API ได้}
    \begin{mycustomenum2}
        \item ตรวจสอบว่า Momoi API Service ทำงานอยู่ตาม endpoint ที่กำหนดในไฟล์ Environment
        \item ตรวจสอบการตั้งค่า API URL และ AUTH API URL ในไฟล์ Environment ให้ชี้มายังบริการ Momoi อย่างสอดคล้องกัน
        \item ตรวจสอบ CORS Settings และ cookie configuration ของบริการปลายทาง
    \end{mycustomenum2}
\end{mycustomenum2}

% \newpage
\vspace{0.5em}
\textbf{การเข้าถึงข้อมูลเพิ่มเติม}
\vspace{0.5em}

\begin{mycustomenum2}
    \item อ่าน Documentation เพิ่มเติมใน README.md ของแต่ละ Component
    \item ตรวจสอบ Log Files สำหรับข้อมูลการ Debug
    \item ใช้ Jaeger UI สำหรับติดตามการทำงานของ API
    \item ใช้ Grafana Dashboard สำหรับติดตามประสิทธิภาพระบบ
\end{mycustomenum2}
